\documentclass[a4paper]{article}

\usepackage{graphicx}
\usepackage{amsmath,amssymb}
\usepackage{minted}
\usepackage{multirow}
\usepackage{float}
\usepackage{algorithm}
\usepackage{algpseudocode}
\usepackage{todonotes}
\usepackage{forest}
\usepackage{xstring}
\usepackage{xcolor}
\usepackage[most]{tcolorbox}
\usepackage{xparse}
\usepackage{natbib}

\usetikzlibrary{graphs,graphdrawing}
\usegdlibrary{layered}

\input{/home/donkarlo/Dropbox/repo/nd_language_ai_project/src/nd_language_ai/formal/computer/programming/tex/latex/code_clips/external_graphviz.tex}
\input{/home/donkarlo/Dropbox/repo/nd_language_ai_project/src/nd_language_ai/formal/computer/programming/tex/latex/parts/control_sequence/kind/semtag/semtag.tex}
\input{/home/donkarlo/Dropbox/repo/nd_language_ai_project/src/nd_language_ai/formal/computer/programming/tex/latex/code_clips/seehere.tex}

\title{Large Source Typing Performance}
\date{2026-09-04}
\author{Mohammad Rahmani}

\begin{document}
   \maketitle

   \section{Problem}
   Large Persian LaTeX sources, including \texttt{evil\_farsi\_reconstructed.tex},
   could become visibly slow while typing. The large-source editor still copied
   and re-split the complete visible Source window after every key press, and the
   Source widget used the rich-text \texttt{QTextEdit} document layout.

   \section{Fix}
   \begin{itemize}
      \item The common Source editor now uses \texttt{QPlainTextEdit} and its
      simplified plain-text document layout.
      \item \texttt{LargeLatexEditor} consumes
      \texttt{QTextDocument.contentsChange} and applies only the changed range
      to the canonical chunk buffer.
      \item Moderate large sources remain fully materialized. Source-window
      virtualization is reserved for files of at least 5,000,000 characters.
      \item Large-source syntax highlighting has an explicit low-cost mode, so
      a 40k window is not incorrectly treated as a small document.
      \item Decorative wrap markers, indentation-guide painting, and hanging
      indentation are disabled in large-document mode while editing semantics,
      RTL direction, bookmarks, Preview, and Visual synchronization remain.
   \end{itemize}

   \section{Regression coverage}
   The regression tests verify the \texttt{QPlainTextEdit} base, incremental
   \texttt{contentsChange} synchronization, large highlighter mode, the 5 MB
   virtualization boundary, and randomized incremental buffer edits against a
   canonical Python string.

   \bibliographystyle{plainnat}
   \bibliography{/home/donkarlo/Dropbox/repo/nd_shared_research/bibliography/bibliography.bib}
\end{document}
